\documentclass[11pt]{article}

\usepackage[margin=1in]{geometry}
\usepackage{amsmath}
\usepackage{booktabs}
\usepackage{hyperref}
\usepackage{graphicx}
\usepackage{listings}
\usepackage{xcolor}
\usepackage{float}
\usepackage{parskip}

\hypersetup{
    colorlinks=true,
    linkcolor=blue,
    urlcolor=blue,
    citecolor=blue
}

\lstset{
    basicstyle=\ttfamily\small,
    backgroundcolor=\color{gray!10},
    frame=single,
    breaklines=true,
    columns=flexible
}

\title{
    \textbf{Conversational AI for E-Commerce} \\
    \large A Vision-Based Approach via CLIP, FAISS, and Claude Sonnet\\[0.5em]
    \normalsize Generative AI: Principles and Applications Final
}
\author{Ethan Chen, Gavin Kunish}
\date{May 2026}

\begin{document}

\maketitle

\begin{abstract}
The report outlines the design and implementation of a multimodal conversational chatbot
for e-commerce product support. It combines CLIP embeddings, a FAISS vector index, and Claude claude-sonnet-4-6 to answer product related questions using both text and image inputs.Built on the Kaggle Amazon Product Dataset 2020 of 10,002 products. The full system is deployed
as a Streamlit application that supports natural language queries, product image uploads, and
three prompting modes (zero-shot, few-shot, and multi-shot).
\end{abstract}

\tableofcontents
\newpage

% -----------------------------------------------------------------------
\section{Introduction}
% -----------------------------------------------------------------------

Customer support in e-commerce is expensive and often bottlenecked by the sheer variety of
product-related questions. A customer might ask about specifications, usage instructions,
compatibility, or simply want to know what a product in an image is. Traditional text-only
chatbots cannot handle the image input case, and retrieval-only systems cannot generate
coherent, context-aware answers.

A multimodal Retrieval-Augmented Generation (RAG) system addresses both of these gaps.
By using a vision-language model for retrieval and a large language model (LLM) for generation,
the system can match a query whether it comes in as text or an image to relevant products in
a catalog and then synthesize a helpful, grounded response.

The project builds such a system end-to-end:
\begin{itemize}
    \item \textbf{Data:} Amazon Product Dataset 2020, 10,002 products
    \item \textbf{Encoder:} CLIP ViT-B/32 for shared text-image embeddings
    \item \textbf{Index:} FAISS \texttt{IndexFlatIP} for exact cosine-similarity search
    \item \textbf{LLM:} claude-sonnet-4-6 via the Anthropic API
    \item \textbf{Interface:} Streamlit web application
\end{itemize}

% -----------------------------------------------------------------------
\section{Related Work}
% -----------------------------------------------------------------------

\subsection{CLIP}
CLIP (Contrastive Language-Image Pre-training)~\cite{radford2021clip} trains a text encoder
and an image encoder jointly on 400 million image-text pairs from the web. The training
objective pulls matching text-image pairs together in a shared embedding space while pushing
non-matching pairs apart. The result is an encoder that can produce semantically comparable
vectors for text and images---which is exactly what is needed for cross-modal retrieval.

\subsection{Retrieval-Augmented Generation}
RAG~\cite{lewis2020rag} combines a retrieval step with a generative model. At inference time,
relevant documents are fetched from an external knowledge base and passed as context to the LLM.
This keeps the LLM from hallucinating product facts it was never trained on, and makes the
knowledge base easy to update without retraining.

\subsection{Large Language Models}
Modern LLMs like Claude can follow complex instructions, maintain multi-turn conversation history,
and process image inputs natively. This makes them well-suited as the generation component in a
multimodal RAG pipeline.

% -----------------------------------------------------------------------
\section{Data}
% -----------------------------------------------------------------------

\subsection{Dataset}
The Amazon Product Dataset 2020 is available on Kaggle. The version used here contains
10,002 products across several categories, with Toys \& Games (66.6\%) being the largest,
followed by Home \& Kitchen (7.1\%) and Clothing, Shoes \& Jewelry (6.3\%).

\subsection{Fields and Preprocessing}
The dataset contains 28 columns. After reviewing them, the following five were chosen for
building product descriptions:

\begin{table}[H]
\centering
\begin{tabular}{ll}
\toprule
\textbf{Field} & \textbf{Usage} \\
\midrule
Product Name     & Primary identifier; placed first in the description \\
Brand Name       & Appended as ``Brand: X'' \\
Category         & Top-level only (text before first pipe character) \\
Selling Price    & Appended as ``Price: X'' \\
About Product    & First 300 characters; pipe delimiters replaced with periods \\
\bottomrule
\end{tabular}
\caption{Product fields used for CLIP text descriptions}
\end{table}

Fields like \texttt{Product Specification} and \texttt{Technical Details} were excluded
because they contain heavy HTML-like markup that degrades CLIP text quality. The five
selected fields cover the information customers most commonly ask about.

The combined description is passed through the CLIP processor with
\texttt{truncation=True, max\_length=77} to respect the model's token limit.
Image URLs are preserved in metadata but are not embedded at index time.

% -----------------------------------------------------------------------
\section{System Architecture}
% -----------------------------------------------------------------------

The pipeline has three stages: encoding, retrieval, and generation.

\begin{enumerate}
    \item \textbf{Encoding.} At index build time, each product description is embedded with
    CLIP's text encoder into a 512-dimensional L2-normalized vector. At query time, the user's
    text or uploaded image is embedded with the corresponding CLIP encoder.
    \item \textbf{Retrieval.} The query vector is compared against all 10,002 stored vectors
    using FAISS inner product search. The top-k most similar products are returned.
    \item \textbf{Generation.} The retrieved product metadata (name, brand, price, description)
    is formatted into a context block and passed to Claude along with the user's question.
    Claude generates a natural language response grounded in that context.
\end{enumerate}

\subsection{CLIP Model}
\texttt{openai/clip-vit-base-patch32} was used. The ViT-B/32 variant processes images as
$14 \times 14$ patches of $32 \times 32$ pixels and produces 512-dimensional embeddings for
both modalities. Text and image embeddings are projected into the same space via
\texttt{text\_projection} and \texttt{visual\_projection} layers respectively.

One implementation challenge was that transformers version 5.5.0 changed the return type of
\texttt{get\_text\_features()} and \texttt{get\_image\_features()}. These methods returned a
\texttt{BaseModelOutputWithPooling} object instead of a tensor. The fix was to call the
sub-models directly:

\begin{lstlisting}[language=Python]
# Text embedding
out   = model.text_model(input_ids=..., attention_mask=...)
feats = model.text_projection(out.pooler_output)
feats = feats / feats.norm(dim=-1, keepdim=True)

# Image embedding
out   = model.vision_model(pixel_values=...)
feats = model.visual_projection(out.pooler_output)
feats = feats / feats.norm(dim=-1, keepdim=True)
\end{lstlisting}

\subsection{FAISS Index}
\texttt{faiss.IndexFlatIP} performs exact inner product search. Since all vectors are
L2-normalized, inner product equals cosine similarity. The exact index is appropriate at
10k scale; for millions of products an approximate index (\texttt{IndexIVFFlat} or HNSW)
would be needed.

\subsection{LLM and Prompting}
Claude claude-sonnet-4-6 was chosen because it supports both text and image in the same API
request, handles multi-turn conversation history natively, and follows system prompts reliably.

Three prompting modes are implemented:

\begin{table}[H]
\centering
\begin{tabular}{lp{9cm}}
\toprule
\textbf{Mode} & \textbf{Description} \\
\midrule
Zero-shot & System prompt defines the role and task only. No examples. \\
Few-shot  & System prompt includes 3 worked Q\&A examples covering features,
            comparisons, and image identification. \\
Multi-shot & System prompt includes 5 examples covering price, age suitability,
             product comparison, usage, and color/variant queries. \\
\bottomrule
\end{tabular}
\caption{Prompt modes and their descriptions}
\end{table}

The user question is always framed as: ``Product context from our database: [context]
Customer question: [question].'' Conversation history is passed on every turn so follow-up
questions work correctly.

% -----------------------------------------------------------------------
\section{Implementation}
% -----------------------------------------------------------------------

\subsection{Index Building (\texttt{build\_index.py})}
This one-time script loads the CSV, constructs product descriptions, encodes them in batches
of 128, and saves the FAISS index and JSON metadata to disk. Total time on CPU is roughly
4--5 minutes for 10,002 products.

\subsection{Retrieval Module (\texttt{rag.py})}
Module-level singletons hold the FAISS index, product metadata, CLIP model, and processor.
They are loaded once on first call to \texttt{retrieve()} and reused for all subsequent calls.
The \texttt{retrieve()} function accepts either a text query or a PIL image.

\subsection{Streamlit Application (\texttt{app.py})}
The app has three tabs:
\begin{itemize}
    \item \textbf{Chat:} Main chatbot interface. Users type questions or upload a product
    image. Retrieved products and their images are displayed below each assistant response.
    \item \textbf{Evaluation:} Shows the pre-computed Recall@1/5/10 metrics and a bar chart.
    A ``Run Evaluation Now'' button re-runs the evaluation live.
    \item \textbf{Documentation:} Architecture overview, data preprocessing details, model
    descriptions, and file structure.
\end{itemize}

\texttt{@st.cache\_resource} is used for the RAG resources (loaded once per Streamlit server
process). \texttt{@st.cache\_data} is used for product image downloads.

% -----------------------------------------------------------------------
\section{Evaluation}
% -----------------------------------------------------------------------

\subsection{Test Set Construction}
200 products are sampled uniformly at random (seed=42) from the index. Each product's
exact name is used as the query. The ground truth is that the queried product should appear
in the top-k retrieved results among the 10,001 competing distractors.

\subsection{Results}

\begin{table}[H]
\centering
\begin{tabular}{lr}
\toprule
\textbf{Metric} & \textbf{Score} \\
\midrule
Accuracy@1  & 0.940 \\
Recall@1    & 0.940 \\
Recall@5    & 0.985 \\
Recall@10   & 0.990 \\
\bottomrule
\end{tabular}
\caption{Retrieval evaluation results (n=200 test queries)}
\end{table}

A Recall@1 of 0.94 means that 94\% of queries have the correct product as the single
top result. Recall@10 of 0.99 means that 99\% of queries have the correct product
somewhere in the first 10 results. These are strong results for exact-name queries on
a 10k-product index.

\subsection{Discussion}
The high recall is expected when the query is the product's exact name, since CLIP encodes
names consistently. The 6\% miss rate at Recall@1 comes from a few failure modes: products
with very generic names (e.g., ``Blue T-Shirt'') that match multiple similar items, and
products with names containing special characters or model numbers that CLIP tokenizes
differently than expected.

\subsection{Limitations of the Evaluation}
This evaluation only covers text-to-text retrieval. Cross-modal retrieval (image query
against text index) is harder to measure without a manually annotated image test set, which
the Amazon Product Dataset does not include. Response quality (whether the LLM's answer is
accurate and helpful) is also not formally measured here and would require human evaluation
or a dedicated benchmark.

% -----------------------------------------------------------------------
\section{Challenges and Solutions}
% -----------------------------------------------------------------------

\textbf{Transformers API change.}
Transformers v5.5.0 changed the return signature of \texttt{get\_text\_features()} and
\texttt{get\_image\_features()}, returning a dataclass instead of a tensor. The fix was to
call \texttt{text\_model} and \texttt{vision\_model} directly and apply the projection layers
manually. This is more explicit and version-stable.

\textbf{Image URL availability.}
Amazon CDN image URLs from 2020 are frequently unavailable or return 403 errors. Rather than
embedding images at index time (which would have failed for most products), image URLs are
stored in metadata and fetched on-demand in the Streamlit app with error handling. This keeps
the index fast to build while still displaying images when they are available.

\textbf{CLIP token limit.}
CLIP's text encoder has a hard limit of 77 tokens. Long product descriptions were silently
truncated by the processor, sometimes cutting off key features. The fix was to pre-truncate
the ``About Product'' field to 300 characters before concatenating, so the most important
parts (name, brand, category, price) always fit within the token budget.

\textbf{Streamlit session state with image uploads.}
Streamlit re-runs the entire script on each interaction. Without care, the uploaded file
reference would reset between the upload and the chat submission. The fix was to read the
file bytes immediately and store them in a local variable, and to call \texttt{seek(0)} after
reading so the file object is still valid for display.

% -----------------------------------------------------------------------
\section{Future Improvements}
% -----------------------------------------------------------------------

\textbf{Fine-tuned CLIP.}
The off-the-shelf CLIP model was not trained on e-commerce data. Fine-tuning on Amazon
product image-description pairs (using the dataset's image URLs) would improve cross-modal
alignment for product-specific visual features like color, material, and form factor.

\textbf{Image embedding at index time.}
Currently, only text embeddings are stored. Downloading and embedding product images would
enable true image-to-image similarity search and improve image-query results.

\textbf{Hybrid retrieval.}
Combining CLIP's semantic search with keyword-based BM25 retrieval would handle queries
like exact model numbers (``B07KMVJJK7'') that CLIP handles poorly because they are
out-of-vocabulary.

\textbf{Approximate indexing.}
At 10k products, brute-force \texttt{IndexFlatIP} is fast enough. For a production catalog
with millions of products, an HNSW or IVF index would reduce query latency significantly.

\textbf{Response evaluation.}
Adding an automated evaluation of LLM response quality using a judge model or RAGAS
framework would give a more complete picture of end-to-end system performance beyond
retrieval recall.

% -----------------------------------------------------------------------
\section{Conclusion}
% -----------------------------------------------------------------------

This project demonstrates that a relatively simple RAG pipeline CLIP embeddings, FAISS
retrieval, and a capable LLM can power a practical multimodal e-commerce chatbot. The
retrieval stage is highly accurate (Recall@10 = 0.99) and the system handles both text and
image queries cleanly. The main limitations are around image URL reliability and the lack
of cross-modal evaluation data, both of which are dataset constraints rather than
architectural ones.

The full source code, FAISS index, and metadata are included in the project submission.
The Streamlit app can be run locally with \texttt{streamlit run app.py} after providing
an Anthropic API key.

% -----------------------------------------------------------------------
\bibliographystyle{plain}
\begin{thebibliography}{9}

\bibitem{radford2021clip}
Radford, A., Kim, J. W., Hallacy, C., Ramesh, A., Goh, G., Agarwal, S., \ldots \& Sutskever, I. (2021).
\textit{Learning transferable visual models from natural language supervision.}
Proceedings of the 38th International Conference on Machine Learning (ICML).

\bibitem{lewis2020rag}
Lewis, P., Perez, E., Piktus, A., Petroni, F., Karpukhin, V., Goyal, N., \ldots \& Kiela, D. (2020).
\textit{Retrieval-augmented generation for knowledge-intensive NLP tasks.}
Advances in Neural Information Processing Systems (NeurIPS).

\bibitem{johnson2019faiss}
Johnson, J., Douze, M., \& Jegou, H. (2019).
\textit{Billion-scale similarity search with GPUs.}
IEEE Transactions on Big Data.

\bibitem{anthropic2024claude}
Anthropic. (2024).
\textit{The Claude 3 Model Family: Opus, Sonnet, Haiku.}
Technical Report.

\end{thebibliography}

\end{document}
